% 4-7-rq-answers.tex
%  Budget: 4pp.  HANDOFF section 7 (RQ text + Claim / Not-a-claim), and the
%  "Wording the answers in 4-7" rules.
%  MANDATE:
%   - answer in the scope the question was asked in
%   - state nulls as findings, in their own sentence
%   - HAC p-values only; bootstrap range wherever the pooled regime result appears
%   - NEVER "beats the state of the art"
%  Rules: decimals in \lr{}; integers Persian.

\section{پاسخ به سؤالات پژوهش}
\label{sec:4-7}

این بخش چهار پرسش پژوهشی فصل نخست را، به‌ترتیب و در همان دامنه‌ای که پرسیده
شده‌اند، پاسخ می‌دهد. هر پاسخ دو بخش دارد: آنچه شواهد پشتیبانی می‌کنند، و آنچه
پشتیبانی نمی‌کنند. بخش دوم به‌اندازه‌ی بخش نخست جدی گرفته شده است، چون فاصله‌ی
میان این دو، همان جایی است که یک نتیجه‌ی درست به ادعایی نادرست بدل می‌شود.

یادآوری یک نکته‌ی صوری پیش از پاسخ‌ها: طرح مصوب، پرسش پژوهشیِ شماره‌گذاری‌شده
ندارد و به‌جای آن شش هدف و شش فرض را برمی‌شمارد. چهار پرسشی که در ادامه پاسخ
داده می‌شود، صورت‌بندیِ عملیاتیِ همان اهداف است، نه افزوده‌ای بر آن‌ها.

\subsection*{پرسش نخست: دقت مقایسه‌ای پنج خانواده‌ی مدلی}

\begin{quote}
پنج مدل منتخب تحت یک پروتکل اعتبارسنجی پیش‌رونده‌ی یکسان چه عملکردی
نسبت به یکدیگر دارند، و کدام تفاوت‌ها از نظر آماری معنادار است؟
\end{quote}

\textbf{پاسخ.} بر ۷۲۸ مبدأ پیش‌بینی از چهارم ژانویه‌ی ۲۰۱۶ تا سی‌ویکم دسامبر
۲۰۱۷، ترتیب \lr{MAE} در افق ساعتی چنین است: \lr{LSTM} با \lr{3.8734}،
\lr{LEAR-LASSO} با \lr{3.8988}، \lr{LightGBM} با \lr{3.9678}، \lr{SARIMAX} با
\lr{4.3507} و مدل ساده با \lr{7.7501}.

اما ترتیب، پاسخ نیست. آزمون دیبولد-ماریانو نشان می‌دهد که سه مدل نخست از یکدیگر
تفکیک‌پذیر نیستند. \lr{LSTM} در برابر \lr{LEAR-LASSO} حتی پیش از هر تصحیحی
\lr{p} برابر \lr{0.4036} می‌دهد. \lr{LSTM} در برابر \lr{LightGBM} برآورد
نقطه‌ای پایین‌تری دارد، اما این تفاوت پس از تصحیح هولم\lr{–}بونفرونی بر
خانواده‌ی ۲۱ آزمونه دوام نمی‌آورد (\lr{p} خام \lr{0.0405}، \lr{p} تصحیح‌شده
\lr{0.1215}).

سه نتیجه‌ی مقاوم باقی می‌ماند. هر چهار مدل آموخته بر مدل ساده برتری دارند با
\lr{p} عملاً برابر صفر. \lr{SARIMAX} از سه مدل دیگر پایین‌تر است، با \lr{p}
تصحیح‌شده‌ی حداکثر \lr{5.1e-04}. و هر دو ترکیب بر تک‌تک اعضای خود برتری دارند،
با \lr{p} تصحیح‌شده‌ی حداکثر \lr{3.4e-05}.

\textbf{تساوی، خودْ یافته است.} اینکه بهترین مدل آماری و بهترین مدل یادگیری
عمیقِ این پژوهش تفکیک‌ناپذیرند، نتیجه‌ای است درباره‌ی اینکه یادگیری عمیق در این
مسئله چه‌قدر می‌افزاید، و نه ناکامی در جداکردن مدل‌ها. طراحیِ برابرِ آزمایش
\lr{—} ماتریس ویژگی یکسان، مبادی یکسان، بودجه‌ی پنجاه‌آزمایشیِ تنظیم برای هر
مدل \lr{—} همان چیزی است که این تساوی را به یافته بدل می‌کند؛ بدون آن،
می‌شد آن را به تنظیمِ ناکافیِ یکی از دو طرف نسبت داد.

\textbf{آنچه ادعا نمی‌شود.} این پاسخ نمی‌گوید یادگیری عمیق به‌طور کلی بر
روش‌های آماری برتری دارد یا ندارد. یک معماری از هر خانواده، یک بازار، یک
پنجره‌ی دوساله، یک بودجه‌ی تنظیم.

\subsection*{پرسش دوم: ترکیب و وزن‌دهی رژیم‌محور}

\begin{quote}
آیا ترکیب وزن‌دار دقت را نسبت به بهترین تک‌مدل بهبود می‌دهد، و آیا
مشروط‌کردن وزن‌ها به رژیم بازار بهبود بیشتری نسبت به وزن‌دهی ایستا ایجاد
می‌کند؟ این بهبود، در صورت وجود، تحت کدام شرایط بازار پدید می‌آید؟
\end{quote}

\textbf{بخش نخست، بله.} ترکیب ایستا با \lr{3.5742} و ترکیب رژیم‌محور با
\lr{3.5569} هر دو از بهترین عضو منفرد \lr{—} \lr{LSTM} با \lr{3.8734} \lr{—}
بهترند، و هر هشت مقایسه‌ی ترکیب با اعضایش پس از تصحیح هولم نیز با \lr{p}
حداکثر \lr{3.4e-05} برقرار می‌ماند. این مقاوم‌ترین یافته‌ی فصل است.

\textbf{بخش دوم، بله اما مشروط.} ترکیب رژیم‌محور بر ترکیب ایستا با \lr{p}
برابر \lr{0.0226} برتری دارد. این آزمون، تنها عضو خانواده‌ی تأییدی است و به
همین دلیل بدون تصحیح چندگانگی در سطح \lr{0.05} خوانده می‌شود؛ ثبت تاریخ‌دارِ آن
در یازدهم ژوئیه‌ی ۲۰۲۶، پیش از آنکه هر ترکیبی بر پنجره‌ی آزمون امتیاز بگیرد، تنها
چیزی است که این خوانش را مجاز می‌کند. دو قید باید همراهش بیاید: اگر همین آزمون
از خانواده‌ی اکتشافی می‌آمد، پس از تصحیح به \lr{0.0903} می‌رسید و معنادار
نمی‌ماند؛ و بوت‌استرپ بلوکی بر کل دوره بازه‌ی \lr{0.0129} تا \lr{0.0571} می‌دهد،
یعنی در بلوک‌های ده‌روزه از آستانه عبور می‌کند.

\textbf{بخش سوم پرسش، پاسخ تیزی دارد و مهم‌ترین بخش این جواب است.} تمام بهبود
از ۷۷ روز پرتنش می‌آید، با \lr{5.5132} در برابر \lr{5.6830} و \lr{p} برابر
\lr{0.0063} \lr{—} که در سراسر بازه‌ی بوت‌استرپ نیز معنادار می‌ماند
(\lr{0.0081} تا \lr{0.0439}). بر ۶۵۱ روز آرام، ترکیب رژیم‌محور به‌اندازه‌ی
ده‌هزارم یورو بدتر است و \lr{p} برابر \lr{0.8465}.

\textbf{نتیجه‌ی تهیِ روزهای آرام، خود یک یافته است.} وزن‌دهیِ رژیم‌محور بهبودی
عمومی نیست؛ سازوکاری است که در شرایط پرتنش کار می‌کند و در شرایط آرام بی‌اثر
است. این دقیقاً همان چیزی است که برایش طراحی شده بود، و تحلیل \lr{SHAP} در
\cref{sec:4-6} از راهی کاملاً مستقل همین را تأیید می‌کند: رفتار مدل در روزهای
پرتنش واقعاً متفاوت است.

\textbf{آنچه ادعا نمی‌شود.} این پاسخ نمی‌گوید وزن‌دهی رژیم‌محور به‌طور کلی بر
وزن‌دهی ایستا برتری دارد. می‌گوید بر این بازار و این پنجره، در روزهای پرتنش
برتری داشت و در روزهای آرام نداشت.

\subsection*{پرسش سوم: تفسیرپذیری}

\begin{quote}
کدام خانواده‌های ویژگی بیشترین سهم را در پیش‌بینی دارند، و آیا این سهم
میان روزهای آرام و پرتنش، و میان افق ساعتی و روزانه، تغییر می‌کند؟
\end{quote}

\textbf{پاسخ.} برای مدل \lr{LightGBM}، پیش‌بینی تولید تجدیدپذیرِ روز هدف
(\lr{exog\_2\_D0}) بزرگ‌ترین سهم منفرد را دارد \lr{—} \lr{5.8422} در حالت آرام
و \lr{5.9761} در حالت پرتنش \lr{—} و در برابر تغییر رژیم پایدار است.

بله، سهم‌ها میان دو رژیم تغییر می‌کنند، و تغییر یکنواخت نیست. اهمیت
\lr{price\_D-1} زیر تنش \lr{76.5} درصد افزایش می‌یابد، از \lr{3.4833} به
\lr{6.1471}، در حالی که پیش‌بینی بار روز هدف (\lr{exog\_1\_D0}) از \lr{3.5197}
به \lr{3.3680} کاهش می‌یابد. مدل در شرایط پرتنش سطح قیمت را از آخرین مشاهده
می‌گیرد و بنیان‌های روز را برای تعیین شکل نگه می‌دارد.

بله، سهم‌ها میان دو افق نیز تغییر می‌کنند. در افق روزانه لگ‌های میانی بیش از
نصف اهمیتشان را از دست می‌دهند (\lr{price\_D-2} از \lr{0.8610} به \lr{0.3778})
در حالی که پیش‌بینی بار روز هدف تقریباً دست‌نخورده می‌ماند (\lr{3.5197} در برابر
\lr{3.4393}) و در حالت آرام به بزرگ‌ترین ویژگی بدل می‌شود.

روز هفته، وقتی لگ‌ها حاضرند، تقریباً بی‌اثر است: \lr{0.1464} در افق ساعتی.

\textbf{یک افشای لازم.} مدلی که در \cref{sec:4-6} توضیح داده شده، برازشی
جداگانه بر پنجره‌ی ۱۰۹۲ روزه‌ی منتهی به سی‌ویکم دسامبر ۲۰۱۷ است و نه همان
مدل‌های اعتبارسنجی پیش‌رونده‌ای که اعداد دقت از آن‌ها گزارش شده؛ توضیح‌دادنِ آن
مدل‌ها روی روزهای آموزششان، تحلیلی درون‌نمونه‌ای می‌بود.

\textbf{آنچه ادعا نمی‌شود.} این‌ها «عوامل قیمت برق آلمان» نیستند و به چهار مدل
دیگر تعمیم نمی‌یابند. \lr{SHAP} خروجی یک مدل را تجزیه می‌کند، نه بازار را.

\subsection*{پرسش چهارم: مسیر مستقیم در برابر مسیر تجمیعی}

\begin{quote}
برای پیش‌بینی میانگین بار پایه‌ی روزانه، کدام مسیر دقت بالاتری دارد، و
آیا این نتیجه برای هر پنج مدل یکسان است؟
\end{quote}

\textbf{پاسخ بخش نخست: تجمیع، برای سه مدل از چهار مدل.} میانگین‌گیری از ۲۴
پیش‌بینی ساعتی برای \lr{LEAR-LASSO} (\lr{2.8389} در برابر \lr{2.8989})،
\lr{LightGBM} (\lr{2.9928} در برابر \lr{3.3012}) و \lr{LSTM} (\lr{2.7804} در
برابر \lr{3.1753}) بهتر عمل می‌کند. مدل ساده در هر دو مسیر یکسان است.

\textbf{پاسخ بخش دوم: خیر، یکسان نیست.} \lr{SARIMAX} استثناست و مسیر مستقیم
برایش بهتر است (\lr{3.1976} در برابر \lr{3.2691}). این استثنا توضیح ساختاری
دارد و در \cref{sec:4-4} آمده است: سری بار پایه‌ی روزانه دقیقاً همان نوع سریِ
هموارِ تک‌بعدی است که مدل‌های کلاسیک سری زمانی برایش ساخته شده‌اند.

اینکه پرسش عامدانه دو بخشی صورت‌بندی شده بود، همین استثنا را به بخشی از پاسخ
بدل می‌کند و نه دردسری برای آن.

\textbf{آنچه ادعا نمی‌شود.} این یک قاعده‌ی کلی درباره‌ی تجمیع زمانی در پیش‌بینی
نیست. چهار مدل، یک بازار، یک پنجره.

\subsection*{یک یافته که هیچ پرسشی آن را نمی‌پوشاند}

آزمون برون‌توزیعی \lr{—} ارزیابی مدل‌های تثبیت‌شده بر داده‌ی زنده‌ی ۲۰۲۶
\lr{—} پس از تصویب طرح پیشنهادی به پروژه افزوده شد و به همین دلیل زیر هیچ‌یک از
چهار پرسش بالا قرار نمی‌گیرد. با این حال، به گمان این پژوهش، مهم‌ترین یافته‌ی
کار است. نتیجه‌ی آن، همراه با آنچه درباره‌ی محدودیت‌های این پژوهش می‌گوید، در
\cref{sec:5-2} گزارش می‌شود.
